\documentclass[10pt,a4paper]{amsart}
\usepackage[margin=19mm]{geometry}
\usepackage{amsmath,amssymb,mathtools}
\usepackage[scr=rsfs,cal=euler]{mathalfa}
\usepackage{xfrac}
\usepackage[unicode,hidelinks,bookmarks=false]{hyperref}
\newcommand{\ie}{i.e.\ }
\newcommand{\cf}{cf.\ }
\newcommand{\resp}{respectively\ }
\newcommand{\Subsection}{\subsection}
\providecommand{\nobreakdash}{\nobreak\hbox{-}}
\setlength{\emergencystretch}{2em}
\multlinegap0pt
\usepackage{mathtools}
\usepackage{amssymb}
\usepackage{tikz}
\usepackage{enumerate}
\newtheorem{prop}{Proposition}
\newtheorem{clist}{Clist}
\newtheorem{thm}{Theorem}
\newtheorem{lemma}{Lemma}
\usepackage{cleveref}
\crefname{prop}{Proposition}{Propositions}
\crefname{clist}{Clist}{Clists}
\crefname{thm}{Theorem}{Theorems}
\crefname{lemma}{Lemma}{Lemmas}
\newcommand{\F}{\mathbb{F}}
\DeclareMathOperator{\GL}{GL}
\DeclareMathOperator{\LL}{\Lambda}
\newcommand{\Z}{\mathbb{Z}}
\renewcommand{\mod}{\,\,\text{\normalfont mod}\,}
\title{Swan modules and homotopy types after a single stabilisation}
\author{Tommy Hofmann, John Nicholson}
\date{Source: arXiv:2507.21975v2 (CC BY 4.0)}
\begin{document}
\maketitle

\noindent\emph{Source and licence.} Swan modules and homotopy types after a single stabilisation. Version arXiv:2507.21975v2 (CC BY 4.0), distributed by arXiv under the Creative Commons Attribution 4.0 International licence, http://creativecommons.org/licenses/by/4.0/, as recorded in the arXiv licence metadata for this exact version and shown on the abstract page https://arxiv.org/abs/2507.21975v2. Full licence notice: \texttt{licenses/arxiv-2507.21975-CC-BY-4.0.txt}.\par\smallskip

\noindent\textit{Editorial scope.} Counted unit: the article's prop prop:swansetup (source0.tex lines 506--532). The article's complete original proof of it is delivered with it (source0.tex lines 831--849, one \texttt{proof} environment of 19 lines, which the article places away from the statement). No mathematics is written, repaired or re-derived: the statement and the proof are the article's own lines, in the article's own notation, and no retained line is substituted or re-flowed.  Retained uncounted as the source-local context the proof needs: lines 598--604; lines 643--645; lines 721--729; lines 815--817.  Nothing was omitted from any retained span, so the package carries no omission record.  No retained line contains a citation, so the wrapper carries no bibliography and no bibliography entry.  No retained line contains a figure environment, an image-inclusion command or drawing code, so no image is delivered and the package declares no asset dependency.  Editorial formatting only, in addition: an AMS \texttt{amsart} preamble that restates the article's own declarations and macros used by the retained lines, namely the environments \texttt{prop}, \texttt{clist}, \texttt{thm}, \texttt{lemma} on separate counters, together with the standard packages those lines need.  The retained environments print in this document as Proposition 1, Clist 1, Theorem 1, Lemma 1 in order of appearance, whereas the article prints the same items with the numbering of its own full document.  The complete native source of the article as distributed in the arXiv e-print for this exact version is bundled unchanged as \texttt{sources/182385/source0.tex}.
\begin{equation}\label{eq:rhos}
\begin{aligned}
\rho_1 &: \F_2 D_{4p}^\times \xrightarrow[]{f_*} \F_2 V^\times \twoheadrightarrow {(\F_2 V)^\times}/V = \Z/2 \cdot (1+x+y) \cong \Z /2 ,
\\
\rho_2 &: \F_2 D_{4p}^\times \to \GL_2(Z/2Z) \xrightarrow[]{\det} (Z/2Z)^\times \xrightarrow[]{\cong} (R_p/2R_p)[C_2]^\times  \twoheadrightarrow \frac{(R_p/2R_p)[C_2]^\times}{1+\mathfrak{p}R_pC_2}.
\end{aligned} 
\end{equation}

\begin{thm}\label{cor:N15}
The Swan module $(N, 15)$ of $\Z[Q_{56}]$ is a non-free stably free $\Z[Q_{56}]$-module.
\end{thm}

\begin{lemma}\label{lemma:rho2u}
The following holds over $Q_{56}$:
\begin{clist}{(i)}
    \item
        $\rho_2(u_{(N, 15)}) \neq 0$.
    \item
        $\rho_2(u_{(N, 9)}) = 0$.
\end{clist}
\end{lemma}

\begin{lemma}\label{lemma:rho2units}
  For $p = 7$ we have $\rho_2(\Lambda^\times) = \{0\}$.
\end{lemma}

\begin{prop}\label{prop:swansetup}
Let $p$ be an odd prime and let $r \in \Z$, $\gcd(r, 8p) = 1$.
Let $\zeta_p$ be a primitive $p$th root of unity, let $\lambda_p = \zeta_p+\zeta_p^{-1}$, let $R_p = \Z[\lambda_p]$ and let $\mathfrak p$ be the unique prime ideal of $R_p$ over $p$.
\begin{clist}{(i)}
\item
    We have $(a_1)_\#((N, r)) \cong \Z D_{4p}$ and $(a_2)_\#((N, r)) \cong \Lambda$.
\item
Let $s,t \in \Z$, $s \geq 1$ be such that $rs = 1+t|G|$ and $m = (s-1)/2$. The isomorphism class of $(N, r)$ is determined by the double coset $[u_{(N, r)}] \in \Z D_{4p}^\times \backslash \F_2 D_{4p}^\times /{\LL^\times}$, where
\[ u_{(N,r)} = 1+(1+x^r + \dotsb +  (x^r)^{m-1})x^{r-1}(x+y) + tN \in \F_2 D_{4p}^\times. \]
\item
There exist maps $\rho_1 \colon \F_2 D_{4p}^\times \to \Z/2$, $\rho_2 \colon \F_2 D_{4p}^\times \to (R_p/2R_p)[C_2]^\times/(1 + \mathfrak p R_p C_2)$ inducing a bijection
\[
\overline \rho \colon \Z D_{4p}^\times \backslash \F_2 D_{4p}^\times /{\LL^\times}
\longrightarrow \left(\Z/2 \oplus \frac{(R_p/2R_p)[C_2]^\times}{1+\mathfrak{p}R_pC_2}\right) \Big/ (\rho_1, \rho_2)(\Lambda^\times), \quad [u] \mapsto [(\rho_1(u),\rho_2(u))]
\] 
with $\overline \rho([1]) = 0$.
In particular, $(N, r)$ is free if and only if $\overline \rho([u_{(N, r)}]) = 0$.
\item
We have
\[ \rho_1(u_{(N,r)}) = 
\begin{cases}
0, & r \equiv \pm 1 \mod 8 \\
1, & r \equiv \pm 3 \mod 8. 
\end{cases}
\]
\end{clist}
\end{prop}

\begin{proof}[Proof of \cref{cor:N15}]
We have $Q_{56} = Q_{8 \cdot 7}$ and hence $p = 7$ with the notation of the preceding sections.
Let $u = u_{(N, 15)}$
and let $X = {(R_p/2R_p[C_2]^\times)}/({1+\mathfrak{p}R_pC_2})$.
Recall that by~\cref{prop:swansetup}~(iii), the maps $\rho_1$ and $\rho_2$ from~\cref{eq:rhos} induce a bijection
\[
\Z D_{4p}^\times \backslash \F_2 D_{4p}^\times /{\LL^\times}
\longrightarrow (\Z/2 \oplus X )/ (\rho_1,\rho_2)(\Lambda^\times).
\]
Now the composition of the canonical projections
$\Z/2 \oplus X \to
X
\to X/\rho_2(\Lambda^\times)
$ is trivial on $(\rho_1,\rho_2)(\Lambda^\times)$.
Thus we have an induced homomorphism
\[ \F_2 D_{4p}^\times \to \Z D_{4p}^\times \backslash \F_2 D_{4p}^\times /{\LL^\times } \to (\Z/2 \oplus X)/(\rho_1, \rho_2)(\Lambda^\times) \to X/\rho_2(\Lambda^\times),\]
which is in fact equal to $\rho_2$ by \cref{lemma:rho2units}.
Since $\rho_2(u_{28, 15}) \neq 0$ by \cref{lemma:rho2u}~(i), we conclude that $[u_{28, 15}]$ is not the trivial double coset and hence $(N, 15)$ not free by \cref{prop:swansetup}~(i).
\end{proof}

\end{document}
